A documentação de engenharia de software foi organizada retrospectivamente nesta revisão, a partir do código existente. Os requisitos da Tabela~\ref{tab:requisitos-pap} sintetizam comportamentos e qualidades pretendidos; não se afirma que constituíam uma especificação formal aprovada antes da implementação. A Figura~\ref{fig:arquitetura-pap} apresenta as responsabilidades dos componentes. Diagramas comportamentais e a rastreabilidade detalhada estão no Apêndice~\ref{ap:engenharia}.
\begin{table}[H]\centering\small
\caption{Síntese dos requisitos da extensão São Paulo.}\label{tab:requisitos-pap}
\begin{tabular}{p{1.5cm}p{11.5cm}}\hline
ID & Comportamento ou qualidade pretendida\\\hline
RF01 & Identificar o embarque pela localização, pelo QR Code ou pela seleção manual.\\
RF02 & Consultar partidas e distinguir previsão de estimativa programada.\\
RF03 & Planejar por destino com caminhada, espera e baldeações, aceitando coordenadas por toque prolongado.\\
RF04 & Comparar alternativas e abrir a rota ou um trecho no mapa.\\
RF05 & Salvar, renomear e remover favoritos locais, confirmando a remoção.\\
RF06 & Cancelar mediante confirmação, preservando embarque e favoritos; permitir novo destino.\\
RF07 & Acompanhar posição do ônibus e previsão de embarque compatível com a caminhada, com atualização a cada 15 segundos.\\
RQ01 & Informar progresso sem perder a seleção ou redefinir o zoom.\\
RQ02 & Reduzir esforço de interação preservando compreensão e controle.\\\hline
\end{tabular}\fonte{Síntese do grupo baseada na implementação e nos objetivos de IHC.}
\end{table}

\begin{figure}[H]
\centering\begin{tikzpicture}[>=Stealth,font=\small,box/.style={draw,rounded corners,align=center,text width=3.5cm,minimum height=1cm,fill=white},arr/.style={<->,thick}]
\draw[rounded corners,fill=blue!4] (-2.2,1.1) rectangle (2.2,-7.7);
\node[font=\bfseries] at (0,.7) {PaP};
\draw[rounded corners,fill=black!3] (4.1,1.1) rectangle (8.5,-7.7);
\node[font=\bfseries,align=center] at (6.3,.55) {Fontes e serviços\\externos};
\node (p) at (0,2) {Passageiro};
\node[box] (ui) at (0,-.6) {Interface web\\React / Leaflet};
\node[box] (local) at (0,-2.6) {Armazenamento do navegador\\Favoritos};
\node[box] (api) at (0,-4.6) {Backend Django / Python\\Consultas e planejamento};
\node[box] (db) at (0,-6.8) {SQLite e arquivos locais\\Índice GTFS e cache};
\node[box] (tiles) at (6.3,-.6) {OpenStreetMap\\Mapa-base};
\node[box] (ext) at (6.3,-4.6) {Olho Vivo, busca de endereços e roteamento a pé};
\node[box] (gtfs) at (6.3,-6.8) {GTFS SPTrans\\Cadastro e programação};
\draw[->,thick] (p)--(ui);
\draw[arr] (ui)--(local);
\draw[arr] (ui.west)--++(-.7,0)|-(api.west);
\node[rotate=90,font=\footnotesize,fill=white] at (-2.55,-2.6) {JSON / NDJSON};
\draw[arr] (api)--(db);
\draw[->,thick] (tiles)--node[above,font=\footnotesize]{imagens}(ui);
\draw[arr] (api)--node[above,font=\footnotesize]{consultas}(ext);
\draw[->,thick] (gtfs)--node[above,font=\footnotesize]{importação}(db);
\end{tikzpicture}

\caption{Visão de arquitetura da extensão São Paulo, inspirada no nível de contêineres do C4.}\label{fig:arquitetura-pap}
\fonte{Elaboração do grupo, documentação da implementação em 7 out. 2026.}
\end{figure}

A interface consome a API do backend e o mapa-base; os favoritos permanecem no navegador. O backend consulta o índice GTFS e serviços externos, sem transferir credenciais operacionais para a interface. As setas resumem dependências lógicas e não descrevem uma topologia de implantação em produção. A integração externa está sujeita a indisponibilidade, que deve ser comunicada sem transformar estimativas programadas em previsões ao vivo.
